% paper_iii_body.tex

\maketitle

\begin{abstract}
We study finite-sample reconstruction in measure-preserving systems.
Given a time series $x_1,\ldots,x_n$ from a system $(X,\mathcal{B},\mu,T)$
and a single observable $h$, we ask whether reconstruction --- the condition
that the delay algebra $\sigma(h, h\circ T, \ldots)$ generates $\mathcal{B}$
modulo $\mu$ --- is detectable from data alone, at what rate, and with what
certificates.

We prove three theorems.
The \emph{Algebra Theorem} shows that the empirical $\sigma$-algebra
approximation error $\hat\delta(L,n)$ concentrates around the true error
$\delta(L)$, and that the elbow stopping rule $\hat{L}^*$ achieves the
minimax-optimal rate $n^{-s/(2s+d)}$ for $\mathrm{H\ddot{o}lder}(s)$ targets
on a $d$-dimensional state space, without oracle knowledge of the mixing rate,
the lag, or the smoothness index.
The \emph{Dynamics Theorem} shows that under uniform separation, the delay
map $\Phi_h^{(L)}$ is bi-Lipschitz and estimated delay vectors certify point
separation at rate $n^{-\beta/(2\beta+d)}$.
The \emph{Conjunction Theorem} shows that for deterministic $T$, algebra
separation and metric separation are the same event --- not merely correlated,
but identical --- and that both witnesses certify this from data under
reconstruction and uniform separation.

A third computable witness emerges from the fibre structure of the delay map.
Under the fibre mixing condition --- that $\nu_L$-almost every large fibre is
balanced --- the collision entropy $H_2(\nu_L) := -\log(\mu\otimes\mu)(R_L)$
diverges if and only if $\delta(L) \to 0$.
The key is a conditional variance identity expressing $\delta(L)$ as an
integral of within-fibre variance over unseparated trajectory pairs; the
identity reverses under fibre mixing but fails when large fibres are
monochromatic --- entirely inside or outside the hardest-to-approximate event.
All three witnesses --- algebraic, geometric, and information-theoretic ---
measure the same failure of separation across delay fibres.
\end{abstract}

\begin{quote}\small
\textbf{2020 Mathematics Subject Classification.}
62G08 (primary); 37A05, 62M10, 28A60, 94A17 (secondary).

\textbf{Keywords.}
reconstruction, delay embedding, $\sigma$-algebra approximation,
empirical witness, minimax rate, stopping rule, collision entropy,
Rényi entropy, fibre mixing, conditional variance, Takens theorem,
measure-preserving system.
\end{quote}

% ============================================================
\section{Introduction}
\label{III:sec:intro}
% ============================================================

The reconstruction theorem~\cite{mahon_paper2} states that the observable
algebra $\mathcal{O}_h = \sigma\{h \circ T^n : n \in \mathbb{Z}\}$ generates
$\mathcal{B}$ modulo $\mu$ if and only if the delay map $\Phi_h$ is a
measure-theoretic embedding --- an exact asymptotic condition that says
nothing about rates, certificates, or what can be verified from a finite
time series.
The results of~\cite{mahon_paper1,mahon_paper2} are taken as given;
this paper asks what finite data can certify about the structure they establish.

The certification question has two natural levels.  At the level of the
general query system~\cite{mahon_paper1}, one can ask: given finite responses
to a directed family of queries, does the empirical evidence certify that the
observable structure is resolving correctly, and can CE-failure be detected
from data?  Section~\ref{III:sec:general} answers this at full generality, with
no temporal or geometric assumptions.  The central objects are the population
separation defect $\delta(\mathcal{G})$ and its empirical U-statistic proxy;
the main results are a population convergence theorem (CE drives the defect to
zero under honest refinement), an estimation theorem (the empirical proxy
concentrates under IID sampling), and a CE-shadow corollary (a persistent
empirical floor is evidence against CE).

The remainder of the paper works in the dynamical specialisation --- a
measure-preserving system $(X, \mathcal{B}, \mu, T)$ with observable
$h \in L^\infty$ and time series $x_1, \ldots, x_n$ --- where the additional
structure of mixing and geometry makes sharper, rate-optimal guarantees
achievable.  Here the joined query algebra $\mathcal{O}_h^{(L)}$ is an honest
refinement scheme by construction, and the separation defect specialises to
$\deltaL = \delta(\mathcal{O}_h^{(L)})$.  Three theorems certify reconstruction
in this regime.

\textbf{The Algebra Theorem} (Section~\ref{III:sec:algebra}).
The elbow stopping rule $\Lhat$ fires when the empirical defect
$\delhat{L}{n}$ stops decreasing.
Under exponential mixing with resolvability constant $C_\lambda > 4$,
$\Lhat$ achieves the minimax-optimal rate $n^{-s/(2s+d)}$ for
$\mathrm{H\ddot{o}lder}(s)$ targets on a $d$-dimensional state space,
without any oracle knowledge of the mixing rate, the lag $L^*(n)$,
or the smoothness index $s$.

\textbf{The Dynamics Theorem} (Section~\ref{III:sec:dynamics}).
When $T \in C^r$ and the uniform separation condition~(US) holds,
the delay map $\PhiL$ is bi-Lipschitz, and the empirical delay vectors
constitute a second witness certifying that points are being separated by
the dynamics.  This kernel witness converges at rate $n^{-\beta/(2\beta+d)}$
and fails precisely when~(US) fails.

\textbf{The Conjunction Theorem} (Section~\ref{III:sec:conjunction}).
For deterministic $T$, the TV separation $\preddisc{L}(x,x') = 1$ if and
only if $x$ and $x'$ are separated by $\OhL$ --- the algebraic identity
$\sigma(\PhiL) = \OhL$ is the Markov bridge.  Under reconstruction and~(US),
both witnesses certify the same object; the theorem names two failure modes
(reconstruction failure collapses the algebra witness; (US)-failure collapses
the kernel witness) that jointly exhaust the ways certification can fail.

The dynamical specialisation inherits its honest-refinement property from lag
growth, freeing the mixing assumption to buy rate results that the general
theory cannot reach.  The discussion section records where each proof requires
that structure and what would be needed to extend beyond it.

% ============================================================
\section{General observational certification}
\label{III:sec:general}
% ============================================================

The certification question can be posed before any temporal or geometric
structure is assumed.  Fix a query system $(\mathcal{E}_n)_n$ with
observable $\sigma$-algebra $\sigma(\mathcal{C})$ and realization measure
$\mu$ satisfying collective exhaustion (CE)~\cite{mahon_paper1}.  For a
finite joined query algebra $\mathcal{G} = \mathcal{E}_{n_1} \vee \cdots
\vee \mathcal{E}_{n_k}$, the \emph{population separation defect} is
\[
  \delta(\mathcal{G})
  \;:=\;
  \sup_{A \in \sigma(\mathcal{C})} \inf_{E \in \mathcal{G}} \mu(A \triangle E).
\]
It vanishes if and only if $\sigma(\mathcal{G}) = \sigma(\mathcal{C})$
modulo $\mu$ — that is, the queried family has achieved complete observational
resolution.  A refinement scheme $(\mathcal{G}_k)$ is \emph{honest} if it is
increasing and $\sigma(\bigcup_k \mathcal{G}_k) = \sigma(\mathcal{C})$ modulo
$\mu$: the queries genuinely exhaust the observable system.  Under an honest
scheme, $\delta(\mathcal{G}_k)$ is non-increasing.

\begin{theorem}[Population theorem]
\label{III:thm:population}
Under CE and any honest refinement scheme $(\mathcal{G}_k)$,
\[
  \delta(\mathcal{G}_k) \;\longrightarrow\; 0.
\]
\end{theorem}

\begin{proof}
CE rules out persistent mass on globally vanishing events~\cite{mahon_paper1}.
The events in $\sigma(\mathcal{C})$ not approximable by $\mathcal{G}_k$ form
a decreasing family whose intersection is $\mu$-null as the scheme exhausts
$\sigma(\mathcal{C})$.  By CE, their $\mu$-measure tends to zero.
\end{proof}

Given $N$ independent query responses $\omega_1, \ldots, \omega_N$ drawn
from $\mu$, let $R_i = (Q_{n_1}(\omega_i), \ldots, Q_{n_k}(\omega_i)) \in
\alpha_{n_1} \times \cdots \times \alpha_{n_k}$ be the joint response vector
of the $i$-th sample.  Define the \emph{empirical separation defect}
\[
  \hat\delta(\mathcal{G}, N)
  \;:=\;
  \frac{1}{N^2} \sum_{i=1}^N \sum_{j=1}^N \mathbf{1}[R_i = R_j],
\]
the empirical collision probability of the joint response.  This is a
second-order U-statistic requiring no geometric structure on $\Omega$.

\begin{theorem}[Estimation theorem]
\label{III:thm:estimation}
Under IID sampling from the observable law of the joined query responses,
$\hat\delta(\mathcal{G}, N) \to \delta(\mathcal{G})$ in probability as
$N \to \infty$.
\end{theorem}

\begin{proof}
$\hat\delta(\mathcal{G}, N)$ is a second-order U-statistic for a bounded
kernel; the result is an instance of the standard U-statistic law of large
numbers~\cite{Hoeffding1948}.
\end{proof}

\begin{corollary}[CE-shadow]
\label{III:cor:ce-shadow}
Suppose $(\mathcal{G}_k)$ is an honest refinement scheme and samples are
drawn IID.  If $\hat\delta(\mathcal{G}_k, N)$ exhibits a persistent floor
--- failing to decrease below some $c > 0$ as $k \to \infty$ and $N \to
\infty$ --- that is evidence against CE: CE is the population condition that
guarantees the floor can be discharged by further refinement.
\end{corollary}

\begin{remark}[Scope of the CE-shadow corollary]
\label{III:rem:ce-shadow-scope}
The corollary is one-directional.  A persistent empirical floor implies
CE-failure; the converse is not claimed.  The population defect $\delta$
is defined inside $\sigma(\mathcal{C})$, while CE also controls escape to
the Stone boundary — non-principal ultrafilters that do not manifest as a
deficit in the defect functional.  A charge family could achieve
$\delta(\mathcal{G}_k) \to 0$ (complete resolution within $\sigma(\mathcal{C})$)
while still failing CE through mass concentrated on non-realised boundary
points~\cite{mahon_paper1}.

The honest-refinement condition is essential: without it, a stalling empirical
defect could be a design artifact (queries not cutting the unresolved atoms)
rather than a population-level admissibility failure.

This corollary gives CE an empirical signature it does not have in the
abstract treatment of~\cite{mahon_paper1}: collective exhaustion is no longer
only an admissibility condition for deriving $\sigma$-additivity, but also a
structural hypothesis with a detectable finite-sample consequence.
\end{remark}

The remainder of this paper establishes what is achievable in the dynamical
specialisation, where the honest-refinement condition is automatic and mixing
buys rate-optimal guarantees; see Section~\ref{III:sec:intro} for the overview.

% ============================================================
\section{Setup and Notation}
\label{III:sec:setup}
% ============================================================

\subsection{The measure-preserving system}

Let $(X, \Bcal, \mu)$ be a probability space where $X$ is a compact smooth
$d$-dimensional Riemannian manifold and $\mu$ is a Borel probability measure
with smooth, everywhere-positive density $\rho : X \to \mathbb{R}_{>0}$
satisfying
\[
  \rho_{\min} \;:=\; \operatorname{essinf}_X \rho \;>\; 0.
\]
Let $T : X \to X$ be a $\mu$-preserving measurable map.
We make no injectivity or invertibility assumption on $T$ beyond what
is stated in each hypothesis.

\subsection{The observable and the delay map}

Fix an observable $h \in \Lp{\infty}(\mu)$ with $\|h\|_{\Lp{\infty}} =: a < \infty$.
For each lag depth $L \geq 0$ define the \emph{delay map}
\[
  \PhiL : X \to \mathbb{R}^{L+1},
  \qquad
  \PhiL(x) \;=\; \bigl(h(x),\, h(Tx),\, \ldots,\, h(T^L x)\bigr).
\]
The \emph{observable $\sigma$-algebra at lag~$L$} is
\[
  \OhL \;:=\; \sigma(\PhiL) \;=\; \sigma\bigl\{h \circ T^k : 0 \leq k \leq L\bigr\}.
\]
The \emph{delay-polynomial class} of degree $D \geq 1$ is
\[
  \AhL \;:=\;
  \bigl\{p \circ \PhiL : p \text{ is a polynomial of degree} \leq D \text{ on }
  \mathbb{R}^{L+1}\bigr\} \;\subseteq\; \Lp{2}(X, \OhL, \mu).
\]
Since $\PhiL(X) \subseteq [-a,a]^{L+1}$, all elements of $\AhL$ are bounded.

\subsection{The approximation error}

The central quantity is the $\sigma$-algebra approximation error
\[
  \deltaL \;:=\; \delta(\OhL) \;=\;
  \sup_{S \in \Bcal} \inf_{E \in \OhL} \mu(S \triangle E),
\]
which measures how much of $\Bcal$ remains outside $\OhL$.
It is non-increasing in $L$, vanishes if and only if reconstruction holds,
and under (G1)--(G2) on a smooth manifold reaches zero by lag $L_0 \leq d-1$~\cite{mahon_paper2}.

\subsection{Fibre structure and collision probability}
\label{III:subsec:fibres}

The delay map $\PhiL$ partitions $X$ into \emph{fibres}: for each
$z \in \mathbb{R}^{L+1}$, define
\[
  F_z \;:=\; (\PhiL)^{-1}(z), \qquad
  p_z \;:=\; \mu(F_z), \qquad
  \mu_z \;:=\; \mu(\,\cdot\mid F_z).
\]
Call $z$ an \emph{$\varepsilon$-large fibre} if $p_z \geq \varepsilon$.
The push-forward $\nuL := (\PhiL)_*\mu$ is the distribution of
delay vectors.

The set of unseparated pairs $\RL := \{(x,x') : \PhiL(x) = \PhiL(x')\}$
decomposes as $\RL = \bigsqcup_z F_z \times F_z$, giving
\begin{equation}
  (\mu \otimes \mu)(\RL)
  \;=\; \int p_z^2\,d\nuL(z).
  \label{III:eq:collision-prob}
\end{equation}
This quantity is the \emph{collision probability} of the delay-vector
distribution: it is the probability that two independent draws from $\mu$
produce the same delay vector.

\subsection{The data model}

We observe a time series $x_1, x_2, \ldots, x_n$ which is an orbit segment of
$T$: $x_i = T^{i-1}x_1$ for some initial condition $x_1 \in X$ drawn from
$\mu$.
The empirical measure is $\mun := \frac{1}{n}\sum_{i=1}^n \delta_{x_i}$.
We do not assume ergodicity or density of the orbit; each result states the
mixing hypothesis it uses.

\subsection{Structural hypotheses}

The following hypotheses are referenced throughout.
They are not assumed globally; each theorem states exactly which subset it uses.

\begin{enumerate}[label=(G\arabic*)]
  \item \emph{No short periodic orbits:}
    $T$ has no periodic orbit of period $\leq L$ in $\operatorname{supp}(\mu)$.
    This ensures the orbit points $x, Tx, \ldots, T^L x$ are distinct
    $\mu$-a.e., which is used in the Sard--Smale argument for full rank of
    $\PhiL$.

  \item \emph{Non-degenerate observable:}
    $h \in C^2(X)$ with $dh \neq 0$~$\mu$-a.e.
    This places $h$ in the Sard--Smale regular-value residual set and
    ensures $\PhiL$ has rank $\min(L+1,d)$ generically.
\end{enumerate}

\begin{enumerate}[label=(US)]
  \item \emph{Uniform separation:}
    $\inf_{x \in X} \sigma_{\min}(D\PhiL|_x) > 0$,
    where $\sigma_{\min}$ denotes the smallest singular value.
    This is a uniform lower Lipschitz bound on $\PhiL$, strictly stronger than
    (G1)--(G2) which give full rank only $\mu$-a.e.
    Condition~(US) is used only in the Dynamics and Conjunction theorems
    (Sections~\ref{III:sec:dynamics} and~\ref{III:sec:conjunction}).
\end{enumerate}

\subsection{The empirical witness}

The empirical analogue of $\deltaL$ is
\[
  \delhat{L}{n}
  \;:=\;
  \sup_{p \in \PD,\, \|p\|_{\Lp{2}(\mu)} \leq 1}
  \bigl\|p - \Ehat[p \mid \PhiL]\bigr\|_{\Lp{2}(\mun)}^2,
\]
where $\PD$ is the class of degree-$D$ polynomials in the delay coordinates
$\PhiL(x)$ (normalised in $\Lp{2}(\mu)$), and $\Ehat[p \mid \PhiL]$ is the
Nadaraya--Watson kernel regression estimator (Definition~\ref{III:def:nw}).

The main result of Section~\ref{III:sec:concentration} is that $\delhat{L}{n}$
concentrates around $\deltaL$ at rate $\varepsilon_n = C\,n^{-2s/(2s+d)}$
under exponential mixing, provided $d > 2s$.

% ============================================================
\section{Bias Bound and Conditional Variance}
\label{III:sec:bias}
% ============================================================

Let $(X, \Bcal, \mu)$ be a probability space and $\mathfrak{m} \subseteq \Bcal$
a sub-$\sigma$-algebra.
Define the \emph{$\sigma$-algebra approximation error}
\begin{equation}
  \delta(\mathfrak{m})
  \;:=\;
  \sup_{S \in \Bcal} \inf_{E \in \mathfrak{m}} \mu(S \triangle E).
  \label{III:eq:delta-def}
\end{equation}
This measures how well $\mathfrak{m}$ approximates $\Bcal$: it is zero if and
only if $\mathfrak{m} = \Bcal$ modulo $\mu$-null sets.

\begin{lemma}[Conditional variance identity]
\label{III:lem:cond-var}
For any $S \in \Bcal$:
\[
  \mathbb{E}\bigl[\mathrm{Var}(\mathbbm{1}_S \mid \OhL)\bigr]
  \;=\;
  \frac{1}{2}
  \int_{\RL}
  \bigl|\mathbbm{1}_S(x) - \mathbbm{1}_S(x')\bigr|^2
  \,d(\mu \otimes \mu)(x,x').
\]
\end{lemma}

\begin{proof}
By disintegration, $(\mu \otimes \mu)|_{\RL} = \int \mu_z \otimes \mu_z\,d\nuL(z)$.
On each fibre $F_z$:
\[
  \tfrac{1}{2}\int_{F_z \times F_z}
    |\mathbbm{1}_S(x) - \mathbbm{1}_S(x')|^2
    \,d(\mu_z \otimes \mu_z)
  \;=\; \mu_z(S)\,\mu_z(S^c).
\]
Integrating with weight $p_z$:
$\tfrac{1}{2}\int_{\RL}|\mathbbm{1}_S - \mathbbm{1}_{S'}|^2\,d(\mu\otimes\mu)
= \int p_z\,\mu_z(S)\mu_z(S^c)\,d\nuL(z)$.
Since $\mathrm{Var}(\mathbbm{1}_S \mid \OhL)(x) = \mu_z(S)\mu_z(S^c)$ on $F_z$,
both sides equal $\int \mu_z(S)\mu_z(S^c)\,d\nuL(z)$.
\end{proof}

\begin{corollary}[Easy direction]
\label{III:cor:easy}
$\deltaL \;\leq\; \tfrac{1}{2}\,(\mu \otimes \mu)(\RL)$.
\end{corollary}

\begin{proof}
In Lemma~\ref{III:lem:cond-var}, bound $|\mathbbm{1}_S(x) - \mathbbm{1}_S(x')|^2 \leq 1$
and take the supremum over $S$ using
$\deltaL = \sup_S \mathbb{E}[\mathrm{Var}(\mathbbm{1}_S \mid \OhL)]$.
\end{proof}

\begin{remark}[Finite analogue of collective exhaustion]
Corollary~\ref{III:cor:easy} and its converse (Corollary~\ref{III:cor:entropy}, proved
under fibre mixing in \S\ref{III:sec:algebra}) together say: algebraic
indistinguishability $\delta(L)$ vanishes if and only if the mass of
indistinguishable pairs $(\mu\otimes\mu)(R_L)$ vanishes.  This is the
finite-sample counterpart of collective exhaustion~\cite{mahon_paper1}.  There, CE
asserts that set-theoretic vanishing must be witnessed at some finite level of
the refinement order.  Here, the analogous principle is that observational
indistinguishability cannot persist globally without manifesting as a
non-vanishing mass of indistinguishable pairs at finite lag~$L$.  The fibre
mixing condition plays an analogous structural role to CE: it is a sufficient
condition under which the global and pairwise witnesses are comparable.
Whether fibre mixing is itself derivable from dynamical conditions --- or is
genuinely irreducible, as CE is proved to be in~\cite{mahon_paper1} --- remains open;
see the open problems in Section~\ref{III:sec:discussion}.
\end{remark}

\begin{lemma}[Bias bound]
\label{III:lem:bias}
For every $f \in \Lp{\infty}(\mu)$,
\[
  \bigl\| f - \mathbb{E}[f \mid \mathfrak{m}] \bigr\|_{\Lp{2}(\mu)}
  \;\leq\;
  2\|f\|_{\Lp{\infty}} \cdot \delta(\mathfrak{m})^{1/2}.
\]
\end{lemma}

\begin{proof}
\emph{Step~1 (Indicators).}
Fix $S \in \Bcal$.
By definition of $\delta(\mathfrak{m})$, there exists $E \in \mathfrak{m}$
with $\mu(S \triangle E) \leq \delta(\mathfrak{m})$.
Since $\mathbbm{1}_E \in \Lp{2}(X, \mathfrak{m}, \mu)$ and
$\mathbb{E}[\mathbbm{1}_S \mid \mathfrak{m}]$ is the $\Lp{2}$-best
$\mathfrak{m}$-measurable approximation to $\mathbbm{1}_S$,
\[
  \|\mathbbm{1}_S - \mathbb{E}[\mathbbm{1}_S \mid \mathfrak{m}]\|_{\Lp{2}}^2
  \;\leq\;
  \|\mathbbm{1}_S - \mathbbm{1}_E\|_{\Lp{2}}^2
  \;=\; \mu(S \triangle E) \;\leq\; \delta(\mathfrak{m}).
\]

\emph{Step~2 (Bounded $f$ via layer cake).}
For $f \in \Lp{\infty}(\mu)$ with $\|f\|_{\Lp{\infty}} \leq M$, the layer
cake representation gives $f = \int_{-M}^{M} \mathbbm{1}_{S_t}\,dt$ where
$S_t = \{x : f(x) > t\}$.
Since $\mathbb{E}[\,\cdot\mid \mathfrak{m}]$ is linear and continuous in $\Lp{2}$,
\[
  f - \mathbb{E}[f \mid \mathfrak{m}]
  \;=\; \int_{-M}^{M}
    \bigl(\mathbbm{1}_{S_t} - \mathbb{E}[\mathbbm{1}_{S_t} \mid \mathfrak{m}]\bigr)\,dt.
\]
Applying Minkowski's inequality for integrals and Step~1:
\[
  \bigl\|f - \mathbb{E}[f \mid \mathfrak{m}]\bigr\|_{\Lp{2}}
  \;\leq\; \int_{-M}^{M}
    \bigl\|\mathbbm{1}_{S_t} - \mathbb{E}[\mathbbm{1}_{S_t} \mid \mathfrak{m}]\bigr\|_{\Lp{2}}\,dt
  \;\leq\; 2M \cdot \delta(\mathfrak{m})^{1/2}. \qedhere
\]
\end{proof}

We now specialise to the delay setting.
Let $(X, \Bcal, \mu, T)$ be a measure-preserving system and $h \in \Lp{\infty}(\mu)$.
Write $\OhL = \sigma(\PhiL)$ for the observable $\sigma$-algebra at lag~$L$, where
$\PhiL(x) = (h(x), h(Tx), \ldots, h(T^L x))$,
and $\AhL$ for the class of degree-$D$ polynomials in the delay
coordinates~$\PhiL(x)$.
Let $\fhat$ denote the empirical risk minimiser over $\AhL$.

\begin{lemma}[Three-term decomposition]
\label{III:lem:three-term}
For any $f \in \Lp{\infty}(\mu)$, any $L \geq 0$, and any $D \geq 1$:
\begin{align}
  \bigl\|f - \fhat\bigr\|_{\Lp{2}(\mu)}
  &\;\leq\;
    \underbrace{2\|f\|_{\Lp{\infty}} \cdot \delta(L)^{1/2}}_{\text{bias}}
  \;+\;
    \underbrace{\inf_{g \in \AhL}\bigl\|
      \mathbb{E}[f \mid \OhL] - g
    \bigr\|_{\Lp{2}(\mu)}}_{\text{approx.}}
  \;+\;
    \underbrace{\bigl\|g_D^* - \fhat\bigr\|_{\Lp{2}(\mu)}}_{\text{statistical}},
  \label{III:eq:three-term}
\end{align}
where $\delta(L) := \delta(\OhL)$ as in~\eqref{eq:delta-def} and
$g_D^* = \operatorname{argmin}_{g \in \AhL} \|\mathbb{E}[f \mid \OhL] - g\|_{\Lp{2}(\mu)}$.
\end{lemma}

\begin{proof}
Triangle inequality applied to the decomposition
$f - \fhat = (f - \mathbb{E}[f \mid \OhL]) + (\mathbb{E}[f \mid \OhL] - g_D^*)
+ (g_D^* - \fhat)$, with the first term bounded by Lemma~\ref{III:lem:bias}.
\end{proof}

The three-term bound controls the forward error.
For the stopping rule to be honest, we also need a reverse bound: when
$\delta(L)$ is large, the estimator $\fhat$ cannot approximate $f$ well,
regardless of $n$.

\begin{lemma}[Observable reverse lower bound]
\label{III:lem:reverse}
There exists a universal constant $c > 0$ such that for any $f \in \Lp{\infty}(\mu)$
with $\|f\|_{\Lp{\infty}} \geq \frac{1}{2}$ and any $L \geq 0$:
\[
  \bigl\|\mathbbm{1}_S - \fhat\bigr\|_{\Lp{2}(\mun)}
  \;\geq\;
  \frac{\delta(L)^{1/2}}{2} - C' \cdot n^{-s/(2s+d)}
\]
with probability tending to~$1$ as $n \to \infty$, where $S \in \Bcal$ is
the set achieving the supremum in~\eqref{eq:delta-def} at lag~$L$ and
$\mun = \frac{1}{n}\sum_{i=1}^n \delta_{x_i}$ is the empirical measure.
\end{lemma}

\begin{proof}
\emph{Step~A (Population CE gap).}
Let $S^* \in \Bcal$ achieve the supremum in $\delta(L) = \sup_S \inf_{E \in \OhL} \mu(S \triangle E)$.
For any $E \in \OhL$:
\[
  \|\mathbbm{1}_{S^*} - \mathbb{E}[\mathbbm{1}_{S^*} \mid \OhL]\|_{\Lp{2}(\mu)}^2
  \;\leq\;
  \|\mathbbm{1}_{S^*} - \mathbbm{1}_E\|_{\Lp{2}(\mu)}^2
  \;=\; \mu(S^* \triangle E) \;\geq\; \delta(L) - \varepsilon
\]
for any $\varepsilon > 0$.
Taking $\varepsilon \to 0$: $\|\mathbbm{1}_{S^*} - \mathbb{E}[\mathbbm{1}_{S^*} \mid \OhL]\|_{\Lp{2}(\mu)}
\geq \delta(L)^{1/2}$.

Since the conditional expectation is the $\Lp{2}$-projection,
$\|\mathbbm{1}_{S^*} - g\|_{\Lp{2}(\mu)} \geq \delta(L)^{1/2}$ for all $g \in \AhL \subseteq \Lp{2}(X, \OhL, \mu)$.

\emph{Step~B (Empirical vs population norm).}
By Hoeffding's inequality and the bounded difference property of $\mun$,
$\|\mathbbm{1}_{S^*} - g\|_{\Lp{2}(\mun)} \geq \|\mathbbm{1}_{S^*} - g\|_{\Lp{2}(\mu)} - O(n^{-1/2})$
with high probability, uniformly over $g \in \AhL$.

\emph{Step~C (Empirical norm vs $\fhat$).}
Since $\fhat$ minimises empirical $\Lp{2}$ error over $\AhL$, and
$\|\mathbbm{1}_{S^*} - g\|_{\Lp{2}(\mun)} \geq \delta(L)^{1/2} - O(n^{-1/2})$
for all $g \in \AhL$, we have
$\|\mathbbm{1}_{S^*} - \fhat\|_{\Lp{2}(\mun)} \geq \frac{\delta(L)^{1/2}}{2} - C' n^{-s/(2s+d)}$
with high probability, absorbing the $O(n^{-1/2})$ into the statistical rate.
\end{proof}

% ============================================================
\section{Concentration of the Empirical Witness}
\label{III:sec:concentration}
% ============================================================

Throughout this section: $(T,h)$ satisfies (G1)--(G2), $X$ is a compact
smooth $d$-manifold, $\mu$ has smooth positive density, and $d > 2s$.

\subsection{Uniform boundedness}
\label{III:subsec:bounded}

Let $\PD$ denote the class of degree-$D$ polynomials in the delay coordinates
$\PhiL(x) = (h(x),\ldots,h(T^L x))$, normalised so $\|p\|_{\Lp{2}(\mu)} \leq 1$.
Since $h \in \Lp{\infty}$ with $\|h\|_{\Lp{\infty}} =: a < \infty$, the delay
coordinates lie in the compact box $[-a,a]^{L+1}$.

\begin{lemma}[Uniform boundedness]
\label{III:lem:bounded}
Define
\[
  \MLD \;:=\;
  \binom{D+L+1}{L+1}^{1/2}
  \cdot (2D+1)^{1/2}
  \cdot (2a)^{-(L+1)/2}
  \cdot \rho_{\min}^{-1/2},
\]
where $\rho_{\min} = \operatorname{essinf}_X \rho > 0$ is the infimum of
the density of $\mu$.
Then for every $p \in \PD$:
\[
  \|p\|_{\Lp{\infty}} \;\leq\; \MLD \cdot \|p\|_{\Lp{2}(\mu)}.
\]
\end{lemma}

\begin{proof}
Rescale to the unit box by setting $y_i = x_i/a$; polynomials of degree
$\leq D$ in $x \in [-a,a]^{L+1}$ correspond to polynomials $\tilde{p}$ of
degree $\leq D$ in $y \in [-1,1]^{L+1}$.
Expand $\tilde{p}$ in the tensor-product Legendre basis
$\{P_\alpha\}_{|\alpha|\leq D}$, where
$P_\alpha(y) = \prod_{i=1}^{L+1} P_{\alpha_i}(y_i)$ and $\|P_k\|_{\Lp{\infty}[-1,1]} = 1$.
Since $\|P_\alpha\|_{\Lp{\infty}} = 1$, Cauchy--Schwarz gives
\[
  |\tilde{p}(y)|
  \;\leq\; \sum_{|\alpha|\leq D} |c_\alpha|
  \;\leq\; N(D,L+1)^{1/2} \Bigl(\sum_{|\alpha|\leq D}|c_\alpha|^2\Bigr)^{1/2},
\]
where $N(D,L+1) = \binom{D+L+1}{L+1}$ counts the basis elements.
From the $\Lp{2}$ orthogonality
$\|P_\alpha\|_{\Lp{2}[-1,1]}^2 = \prod_i (2\alpha_i+1)^{-1}$,
the coefficient norm satisfies
\[
  \sum_{|\alpha|\leq D}|c_\alpha|^2
  \;\leq\; (2D+1)\,\|\tilde{p}\|_{\Lp{2}}^2,
\]
since the maximum of $\prod_i(2\alpha_i+1)$ over $|\alpha|\leq D$ is
$2D+1$, achieved at $\alpha=(D,0,\ldots,0)$ — the constraint $|\alpha|\leq D$
forces all weight onto one coordinate, giving an $L$-independent maximum.
Converting back via $\|\tilde{p}\|_{\Lp{2}[-1,1]^{L+1}}
= (2a)^{-(L+1)/2}\|p\|_{\Lp{2}(\mathrm{vol})}$ and
$\|p\|_{\Lp{2}(\mathrm{vol})} \leq \rho_{\min}^{-1/2}\|p\|_{\Lp{2}(\mu)}$
yields the result.
The subsequent arguments require $\|h\|_{\Lp{\infty}} \geq \frac{1}{2}$,
equivalently $2a \geq 1$, so that $\MLD$ grows at most polynomially in $L$.
\end{proof}

\subsection{Concentration around the mean}
\label{III:subsec:mcdiarmiid}

Define the function class
\[
  \mathcal{F}_{L,D}
  \;=\;
  \Bigl\{x \mapsto \bigl(p(x) - \Ehat[p \mid \PhiL](x)\bigr)^2
    \;:\; p \in \PD,\; \|p\|_{\Lp{2}(\mu)} \leq 1\Bigr\},
\]
where $\Ehat[p \mid \PhiL]$ is the Nadaraya--Watson estimator
(Definition~\ref{III:def:nw}).
Then $\delhat{L}{n} = \sup_{f \in \mathcal{F}_{L,D}} \frac{1}{n}\sum_{i=1}^n f(x_i)$.

\begin{proposition}[Concentration, {\upshape\textbf{$\star$}}]
\label{III:prop:concentration}
Assume $d > 2s$.
There exist constants $c, C > 0$ depending only on $(d, s, \|h\|_{\Lp{\infty}})$
such that for all $t > 0$ and all $n \geq 1$:
\[
  P\!\bigl(|\delhat{L}{n} - \mathbb{E}[\delhat{L}{n}]| > t\bigr)
  \;\leq\;
  2\exp\!\Bigl(-\frac{c\,n\,t^2}{\MLD^4 \cdot D^d}\Bigr).
  \tag{$\star$}
\]
In particular, setting $t = \varepsilon_n/2$ with
$\varepsilon_n = C'\,n^{-2s/(2s+d)}$:
\[
  P\!\bigl(|\delhat{L}{n} - \mathbb{E}[\delhat{L}{n}]| > \varepsilon_n/2\bigr)
  \;\leq\;
  2\exp\!\Bigl(-c'\,n^{(d-2s)/(2s+d)}\Bigr)
  \;\to\; 0.
\]
\end{proposition}

\begin{proof}
\emph{Step~1 (Boundedness).}
By Lemma~\ref{III:lem:bounded}, every $p \in \PD$ with $\|p\|_{\Lp{2}(\mu)} \leq 1$
satisfies $\|p\|_{\Lp{\infty}} \leq \MLD$.
Hence every $f = (p - \Ehat[p|\PhiL])^2 \in \mathcal{F}_{L,D}$ satisfies
$\|f\|_{\Lp{\infty}} \leq 4\MLD^2 =: B_{L,D}^2$.

\emph{Step~2 (McDiarmid).}
Replacing one observation $x_i$ with $x_i'$ changes
$\delhat{L}{n}$ by at most $B_{L,D}^2/n$.
By the McDiarmid bounded-difference inequality:
\[
  P\!\bigl(|\delhat{L}{n} - \mathbb{E}[\delhat{L}{n}]| > t\bigr)
  \;\leq\; 2\exp\!\Bigl(-\frac{2n^2 t^2}{n \cdot (B_{L,D}^2/n)^2}\Bigr)
  \;=\; 2\exp\!\Bigl(-\frac{n\,t^2}{2B_{L,D}^4/n}\Bigr).
\]

\emph{Step~3 (Rademacher bias).}
By the symmetrisation inequality \cite{BartlettMendelson2002}:
\[
  \mathbb{E}[\delhat{L}{n}] - \delta(L)
  \;\leq\; 2\,\Rad_n(\mathcal{F}_{L,D}).
\]
By the Ledoux--Talagrand contraction lemma, since $f = (p - \cdot)^2$ is
$2B_{L,D}$-Lipschitz in the second argument:
\[
  \Rad_n(\mathcal{F}_{L,D})
  \;\leq\; 2B_{L,D}\,\Rad_n(\PD).
\]
In the Takens regime the image $\PhiL(X)$ is a $d$-dimensional submanifold,
so the VC dimension of $\PD$ restricted to $\PhiL(X)$ is
$O(D^d)$ \cite{Gyorfi2002}.
Hence $\Rad_n(\PD) \leq C\sqrt{D^d/n}$, giving
$\Rad_n(\mathcal{F}_{L,D}) \leq C\,\MLD\sqrt{D^d/n}$.

\emph{Step~4 (Combined).}
Combining Steps~2--3 and absorbing constants:
\[
  P\!\bigl(|\delhat{L}{n} - \delta(L)| > t\bigr)
  \;\leq\; 2\exp\!\Bigl(-\frac{c\,n\,t^2}{\MLD^4\,D^d}\Bigr).
\]
Setting $t = \varepsilon_n/2 = C'n^{-2s/(2s+d)}/2$, the exponent becomes
\[
  c\,n\,\varepsilon_n^2 / (\MLD^4\,D^d)
  \;\propto\; n^{1-4s/(2s+d)}
  \;=\; n^{(d-2s)/(2s+d)}.
\]
This tends to $+\infty$ if and only if $d > 2s$,
proving the final claim.
\end{proof}

\begin{remark}[The $d > 2s$ hypothesis]
\label{III:rem:d>2s}
When $d \leq 2s$, the function class $\mathcal{F}_{L,D}$ is too rich for
$n$ samples to resolve the fine structure of $\deltaL$ at the required scale.
The fix via localised Rademacher complexity \cite{BartlettBousquetMendelson2005}
is technically available but left to future work; we take $d > 2s$ as
a standing hypothesis.
\end{remark}

\subsection{Source~1: the LLN gap}
\label{III:subsec:source1}

\begin{lemma}[LLN gap]
\label{III:lem:source1}
With notation as above:
\[
  \sup_{p \in \PD}
  \Bigl|
    \|p - \mathbb{E}[p \mid \PhiL]\|_{\Lp{2}(\mun)}^2
    - \|p - \mathbb{E}[p \mid \PhiL]\|_{\Lp{2}(\mu)}^2
  \Bigr|
  \;\leq\; C\sqrt{D^d/n}
\]
almost surely as $n \to \infty$, where $C$ depends on $(d,D,\MLD)$.
\end{lemma}

\begin{proof}
The class $\{(p - \mathbb{E}[p|\PhiL])^2 : p \in \PD\}$ is a VC class
of dimension $O(D^d)$.
The uniform Glivenko--Cantelli theorem \cite{Gyorfi2002} gives a.s.\ convergence
of the empirical mean to the population mean, at rate $O(\sqrt{D^d/n})$.
\end{proof}

\subsection{Source~2: estimator bias}
\label{III:subsec:source2}

The estimator $\Ehat[p \mid \PhiL]$ is the Nadaraya--Watson (NW) kernel
regression estimator.

\begin{definition}[Nadaraya--Watson estimator]
\label{III:def:nw}
For bandwidth $h_n > 0$ and kernel $K : \mathbb{R}^{L+1} \to \mathbb{R}_{\geq 0}$
(compactly supported, symmetric, $\int K = 1$):
\[
  \Ehat[p \mid \PhiL(x)]
  \;:=\;
  \frac{\sum_{j=1}^n K\!\bigl(\tfrac{\PhiL(x)-\PhiL(x_j)}{h_n}\bigr)\,p(x_j)}%
       {\sum_{j=1}^n K\!\bigl(\tfrac{\PhiL(x)-\PhiL(x_j)}{h_n}\bigr)}.
\]
\end{definition}

\begin{proposition}[Source~2 estimator bias]
\label{III:prop:source2}
Assume:
\begin{enumerate}[label=(\roman*)]
  \item $T, h \in C^r$ with $r \geq 2$, $\mu$ smooth positive density $\rho$
        with $\rho_{\min} > 0$.
  \item $\PhiL$ has full rank $\mu$-a.e.\ (follows from (G1)--(G2) and
        Lemma~\ref{III:lem:dimeff}).
  \item $T$ is exponentially mixing with rate $\lambda > 0$ and
        $\|h\|_{\Lp{\infty}} \geq \frac{1}{2}$.
  \item The NW bandwidth satisfies $h_n \sim n^{-1/(2\beta+\deff(L))}$
        with $\deff(L) = \min(L+1, d)$ and $\beta \geq 1$.
\end{enumerate}
Then there exists $C_2(D,L) > 0$ such that for all $p \in \PD$:
\[
  \mathbb{E}\|\Ehat[p \mid \PhiL] - \mathbb{E}[p \mid \PhiL]\|_{\Lp{2}(\mu)}^2
  \;\leq\; C_2(D,L)\cdot n^{-2\beta/(2\beta+\deff(L))},
\]
uniformly in $p \in \PD$, where $C_2(D,L) = 2\MLD^2 C_1^2 A^{2\beta}$ and
$A$ is a constant depending on $(\beta, \deff(L), \rho_{\min})$.
\end{proposition}

\begin{proof}
The proof is a four-link chain.

\emph{Link~1 (Smooth disintegration).}
Under conditions~(i)--(ii), the Rokhlin disintegration
$\mu = \int \mu_z\,d((\PhiL)_*\mu)(z)$ has fibre measures $\mu_z$ supported on
the smooth fibres $F_z = \{\PhiL = z\}$.
An implicit-function-theorem argument (three sub-steps: transfer to common
surface via a $C^{r-1}$ diffeomorphism; bound weight variation; bound
Jacobian factor) gives:
\[
  \|\mu_z - \mu_{z'}\|_{\mathrm{TV}} \;\leq\; C_1\,|z - z'|^\beta,
  \qquad \beta = 1 \text{ (Lipschitz floor)},
\]
where $C_1$ depends on $(\|T\|_{C^r}, \|h\|_{C^r}, \rho_{\min})$.
(The optimal exponent $\beta = r-1$ follows from the full Taylor expansion of
the fibre weight; the floor $\beta = 1$ is sufficient for convergence.)

\emph{Link~2 (Bernstein--Markov).}
From Link~1 and Lemma~\ref{III:lem:bounded}: for $p \in \PD$ with
$\|p\|_{\Lp{2}(\mu)} \leq 1$,
\[
  \bigl|\mathbb{E}[p \mid \PhiL = z] - \mathbb{E}[p \mid \PhiL = z']\bigr|
  \;=\; \Bigl|\int p\,d\mu_z - \int p\,d\mu_{z'}\Bigr|
  \;\leq\; \MLD \cdot C_1 \cdot |z - z'|^\beta.
\]
Hence $z \mapsto \mathbb{E}[p \mid \PhiL = z]$ is $\mathrm{H\ddot{o}lder}(\beta)$
with constant $\Lambda_{D,L} := \MLD C_1$, uniformly over $p \in \PD$.

\emph{Link~3 (NW bias-variance balance).}
Standard NW theory \cite{Gyorfi2002} for a $\mathrm{H\ddot{o}lder}(\beta)$
regression function on a $\deff(L)$-dimensional domain gives:
\[
  \mathbb{E}\|\Ehat[p|\PhiL] - \mathbb{E}[p|\PhiL]\|_{\Lp{2}(\mu)}^2
  \;\leq\; \Lambda_{D,L}^2\,h_n^{2\beta}
    + \frac{C_K\,\MLD^2}{n\,h_n^{\deff(L)}\,f_{\min}},
\]
where $C_K$ depends on the kernel $K$ and $f_{\min}$ is the minimum of the
marginal density of $(\PhiL)_*\mu$.
The $\MLD^2$ factor appears in both terms.
Balancing: $\Lambda_{D,L}^2 h_n^{2\beta} = C_K \MLD^2 / (n h_n^{\deff(L)} f_{\min})$
gives $h_n^* \sim A \cdot n^{-1/(2\beta+\deff(L))}$ where $A$ depends only
on $(\beta, \deff(L), C_K, f_{\min}, C_1)$ --- not on $D$, $L$, or $p$.
The $\MLD^2$ factors cancel from both sides of the balance equation.

\emph{Link~4 (Uniformity).}
Substituting $h_n^*$, both terms contribute
$C_2(D,L) \cdot n^{-2\beta/(2\beta+\deff(L))}$ with
$C_2(D,L) = 2\MLD^2 C_1^2 A^{2\beta}$.
Since the bound is $p$-free, it holds uniformly over $\PD$.
\end{proof}

\begin{theorem}[Full convergence $\delhat{L}{n} \to \deltaL$]
\label{III:thm:convergence}
Under the hypotheses of Proposition~\ref{III:prop:source2} and with $d > 2s$:
\[
  \mathbb{E}\bigl|\delhat{L}{n} - \deltaL\bigr|
  \;\leq\;
  C\sqrt{D^d/n}
  \;+\; C_2(D,L)^{1/2} \cdot n^{-\beta/(2\beta+\deff(L))}
\]
and $P(|\delhat{L}{n} - \deltaL| > \varepsilon_n) \to 0$ exponentially,
where $\varepsilon_n = C' n^{-2s/(2s+d)}$.
\end{theorem}

\begin{proof}
Decompose:
\[
  |\delhat{L}{n} - \deltaL|
  \;\leq\;
  \underbrace{|\delhat{L}{n} - \mathbb{E}[\delhat{L}{n}]|}_{\text{concentration, }(\star)}
  \;+\;
  \underbrace{|\mathbb{E}[\delhat{L}{n}] - \deltaL|}_{\text{bias}}
\]
The bias splits as $\text{Source~1} + \text{Source~2}$:
Source~1 is bounded by $C\sqrt{D^d/n}$ (Lemma~\ref{III:lem:source1});
Source~2 is bounded by $C_2(D,L)^{1/2} \cdot n^{-\beta/(2\beta+\deff(L))}$
(Proposition~\ref{III:prop:source2}).
The concentration term is bounded by $(\star)$ (Proposition~\ref{III:prop:concentration}).
Combining at $t = \varepsilon_n/2$ gives the probability statement.
\end{proof}

% ============================================================
\section{The Algebra Theorem}
\label{III:sec:algebra}
% ============================================================

This section assembles the results of Sections~\ref{III:sec:bias}
and~\ref{III:sec:concentration} into the Algebra Theorem: the data-driven stopping
rule $\Lhat$ achieves the minimax-optimal rate $n^{-s/(2s+d)}$ without knowing
the oracle lag $L^*(n)$ or the mixing rate of $T$.

\subsection{Effective dimension and the piecewise rate}

\begin{definition}[Effective dimension]
\label{III:def:deff}
The \emph{effective dimension at lag~$L$} is
\[
  \deff(L) \;:=\; \operatorname{rank}(D\PhiL|_x)
  \quad\text{at a generic point } x \in X.
\]
(The rank is constant on an open dense set by the rank theorem, so this is
well-defined $\mu$-a.e.)
\end{definition}

\begin{lemma}[$\deff$ step function]
\label{III:lem:dimeff}
Under~(G1)--(G2) with $X$ a compact smooth $d$-manifold and $\mu$ smooth
positive density:
\[
  \deff(L) \;=\; \min(L+1,\, d) \quad \text{for all } L \geq 0.
\]
\end{lemma}

\begin{proof}
The upper bound $\deff(L) \leq \min(L+1,d)$ holds because $D\PhiL|_x$ is a
linear map from $T_x X$ (dimension~$d$) to $\mathbb{R}^{L+1}$.

For the lower bound we use the Sard--Smale theorem
(Smale \cite{Takens1981} framework) applied to two failure maps.

\emph{Case $L+1 \leq d$.}
Rank $< L+1$ at $x$ means the cotangent vectors
$\omega_j(x) := (DT^j|_x)^*(dh|_{T^j x})$ are linearly dependent.
Define the failure map
\[
  \Psi : C^2(X) \times X \times \mathbb{P}^L \to T^*X,
  \quad
  (h, x, [\lambda]) \;\mapsto\; d\!\left(\textstyle\sum_j \lambda_j\, h \circ T^j\right)\bigl|_x.
\]
By~(G1), the orbit points $T^j x$ are distinct $\mu$-a.e., so bump functions
at $T^j x$ can prescribe any cotangent vector at $x$ independently.
Hence $D_h\Psi$ is surjective onto $T^*_x X$ (dimension~$d$), and
$\Psi^{-1}(0)$ is a Banach submanifold of codimension~$d$.
The projection to $C^2(X)$ is Fredholm of index $d + L - d = L$, so by
Sard--Smale the regular values form a residual set $R \subset C^2(X)$.
For $h \in R$, the failure locus in $X$ has dimension $\leq L < d$, hence
$\mu$-measure zero.

By~(G2), $dh \neq 0$~$\mu$-a.e., which forces every critical point of
$F_\lambda = \sum_j \lambda_j h \circ T^j$ to lie in the codimension-$d$
submanifold $S_\lambda = \{dh_x \in \operatorname{span}\{\omega_1(x),\ldots,\omega_L(x)\}\}$
(a finite set for each $[\lambda]$, since $X$ is $d$-dimensional).
Hence $h$ is a regular value of $\pi$, i.e., $h \in R$.

\emph{Case $L+1 > d$.}
Immersion failure requires $\exists\, v \neq 0$ with $D\PhiL|_x \cdot v = 0$.
Define
\[
  \Xi : C^2(X) \times (TX \setminus 0) \to \mathbb{R}^{L+1},
  \quad
  (h,(x,v)) \;\mapsto\; \bigl(dh_{T^j x}(DT^j|_x \cdot v)\bigr)_{j=0}^L.
\]
By~(G1), $D_h\Xi$ is surjective onto $\mathbb{R}^{L+1}$, so $\Xi^{-1}(0)$
has codimension $L+1$ and the projection to $C^2(X)$ is Fredholm of index
$2d - 1 - (L+1) = 2d - L - 2$.
For $L \geq d$: $2d - L - 2 \leq d - 2 < d-1$, so the failure locus in $X$
has dimension $< d-1$ and $\mu$-measure zero.
Condition~(G2) places $h$ in the regular-value set by the same argument as
Case~1.
\end{proof}

\begin{proposition}[Piecewise pre-Takens rate]
\label{III:prop:piecewise-rate}
Assume~(G1)--(G2), $X$ compact smooth $d$-manifold, $\mu$ smooth positive
density, $f \in \mathrm{H\ddot{o}lder}(s) \cap \Lp{\infty}(\mu)$, $d > 2s$,
and $T \in C^r$ with $r \geq 2$.
Let $\fhat$ be the empirical risk minimiser over $\AhL$ with $D = D^*(n,L)$
chosen to balance approximation and statistical terms.
Then for each $L \geq 0$:
\[
  \bigl\|f - \fhat\bigr\|_{\Lp{2}(\mu)}
  \;\leq\;
  2\|f\|_{\Lp{\infty}} \cdot \deltaL^{1/2}
  \;+\;
  C(f,d,s) \cdot n^{-s/(2s+\deff(L))}
\]
with probability tending to~$1$, where $\deff(L) = \min(L+1,d)$.
In the Takens regime ($L \geq d-1$): $\deff(L) = d$ and the rate locks at
$n^{-s/(2s+d)}$.
\end{proposition}

\begin{proof}
Apply the three-term decomposition (Lemma~\ref{III:lem:three-term}).
The bias term is bounded by $2\|f\|_{\Lp{\infty}} \cdot \deltaL^{1/2}$
(Lemma~\ref{III:lem:bias}).
By Lemma~\ref{III:lem:dimeff}, $\PhiL(X)$ is a $\deff(L)$-dimensional smooth
submanifold.
On this submanifold, $\mathbb{E}[f \mid \OhL](x) = g^*(\PhiL(x))$ for some
$g^* \in \mathrm{H\ddot{o}lder}(s)$; standard polynomial approximation on a
$\deff(L)$-dimensional smooth domain gives approximation error
$\leq C_{\mathrm{approx}} \cdot D^{-s/\deff(L)}$.
The statistical term is $\leq C_{\mathrm{stat}}(D^{\deff(L)}/n)^{1/2}$ by
Proposition~\ref{III:prop:concentration} with $d$ replaced by $\deff(L)$.
Balancing: $D^* \sim n^{\deff(L)/(2s+\deff(L))}$ and both terms are
$O(n^{-s/(2s+\deff(L))})$.
\end{proof}

\subsection{The oracle lag and the elbow stopping rule}

\begin{definition}[Oracle lag]
\label{III:def:oracle-lag}
The \emph{oracle lag} is
\[
  L^*(n) \;:=\; \min\bigl\{L \geq 0 : \deltaL \leq C_0 \cdot n^{-2s/(2s+d)}\bigr\},
\]
where $C_0 > 0$ is a constant depending on $\|f\|_{\Lp{\infty}}$.
At $L = L^*(n)$, the bias term $2\|f\|_{\Lp{\infty}} \cdot \delta(L^*)^{1/2}$
is $O(n^{-s/(2s+d)})$, matching the statistical floor.
\end{definition}

The oracle lag $L^*(n)$ requires knowledge of the mixing rate, which is
unobservable.
The empirical witness $\delhat{L}{n}$ sidesteps this: it measures reconstruction
quality directly from data.

\begin{definition}[Elbow stopping rule]
\label{III:def:elbow}
Fix a tolerance sequence $\tau_n > 0$.
The \emph{elbow stopping rule} is
\[
  \Lhat \;:=\; \min\bigl\{L \geq 0 :
    \delhat{L+1}{n} - \delhat{L}{n} > -\tau_n\bigr\}.
\]
\end{definition}

\begin{theorem}[Elbow location]
\label{III:thm:elbow}
Assume the hypotheses of Proposition~\ref{III:prop:piecewise-rate} and
exponential mixing $\deltaL \leq A e^{-\lambda L}$ with resolvability constant
$C_\lambda := (1-e^{-\lambda}) A e^\lambda / C' > 4$.
Set $\tau_n = \frac{C_\lambda}{2} \cdot C' \cdot n^{-2s/(2s+d)}$.
Then
\[
  P\!\bigl(\Lhat = L^*(n)\bigr) \;\to\; 1
  \quad\text{exponentially in } n.
\]
\end{theorem}

\begin{proof}
We show the stopping criterion fires at $L = L^*(n)$ (no overshoot)
and not before (no early stop).

\emph{No overshoot: $\Lhat \leq L^*(n)$.}
At $L = L^*(n)$, the population increment satisfies
$|\delta(L^*(n)+1) - \delta(L^*(n))| \leq \delta(L^*(n)) \leq C_0 n^{-2s/(2s+d)}$ (non-increasing $\delta$).
By Theorem~\ref{III:thm:convergence}, $|\delhat{L}{n} - \deltaL| \leq \varepsilon_n = C' n^{-2s/(2s+d)}$ whp.
The empirical increment satisfies
$\delhat{L^*+1}{n} - \delhat{L^*}{n} \geq -C_0 n^{-2s/(2s+d)} - 2\varepsilon_n > -\tau_n$ whp,
so the criterion fires.

\emph{No early stop: $\Lhat \geq L^*(n)$.}
For $L < L^*(n)$, exponential mixing gives population decrement
$\delta(L) - \delta(L+1) \geq \Delta(L) \geq C_\lambda \cdot C' n^{-2s/(2s+d)}$.
Tracking via Theorem~\ref{III:thm:convergence}:
\[
  \delhat{L}{n} - \delhat{L+1}{n}
  \;\geq\; C_\lambda C' n^{-2s/(2s+d)} - 2\varepsilon_n
  \;=\; (C_\lambda - 2)C' n^{-2s/(2s+d)}.
\]
Since $C_\lambda > 4$: $(C_\lambda - 2)C' > \tau_n \cdot 2/C_\lambda \cdot C_\lambda = 2\tau_n/C_\lambda \cdot C_\lambda$,
and specifically $(C_\lambda - 2)C' n^{-2s/(2s+d)} \geq \tau_n$ when
$C_\lambda - 2 \geq C_\lambda/2$, i.e., $C_\lambda \geq 4$.
So the stopping criterion does \emph{not} fire for $L < L^*(n)$, and $\Lhat \geq L^*(n)$.

\begin{remark}[Scope of the elbow theorem]
The resolvability condition $C_\lambda > 4$ is not an artefact: when
$C_\lambda \leq 4$ the elbow drop at $L^*(n)$ merges with the noise floor
$\varepsilon_n$ and the two are indistinguishable from data.
For polynomial mixing, $\Delta(L^*(n)) \sim \varepsilon_n / L^*(n) \to 0$
relative to $\varepsilon_n$, so the elbow drops vanish asymptotically and
the theorem fails.
A separate argument (e.g.\ integrated elbow) is needed for that case and is
not proved here.
\end{remark}
\end{proof}

\begin{theorem}[Algebra Theorem]
\label{III:thm:algebra}
Assume:
\begin{enumerate}[label=(\roman*)]
  \item Hypotheses of Proposition~\ref{III:prop:piecewise-rate}:
    (G1)--(G2), $X$ compact smooth $d$-manifold, $\mu$ smooth positive density,
    $f \in \mathrm{H\ddot{o}lder}(s) \cap \Lp{\infty}(\mu)$, $d > 2s$, $T \in C^r$, $r \geq 2$.
  \item Exponential mixing: $\deltaL \leq A e^{-\lambda L}$ for some $A, \lambda > 0$.
  \item Resolvability: $C_\lambda > 4$.
  \item Exponential mixing for the estimator bias (Proposition~\ref{III:prop:source2})
    and $\|h\|_{\Lp{\infty}} \geq \frac{1}{2}$.
\end{enumerate}
Let $\Lhat$ be the elbow stopping rule (Definition~\ref{III:def:elbow}) with
$\tau_n = \frac{C_\lambda}{2} C' n^{-2s/(2s+d)}$.
Then
\[
  \bigl\|f - \hat{f}_n^{(\Lhat)}\bigr\|_{\Lp{2}(\mu)}
  \;=\; O_p\!\bigl(n^{-s/(2s+d)}\bigr).
\]
The rate $n^{-s/(2s+d)}$ is minimax-optimal for $\mathrm{H\ddot{o}lder}(s)$
functions on a $d$-dimensional domain.
No oracle input (mixing rate, $L^*(n)$, Sobolev index $s$) is required.
\end{theorem}

\begin{proof}
On the event $\{\Lhat = L^*(n)\}$, which has probability tending to~$1$ by
Theorem~\ref{III:thm:elbow}:
\begin{align*}
  \bigl\|f - \hat{f}_n^{(\Lhat)}\bigr\|_{\Lp{2}(\mu)}
  &= \bigl\|f - \hat{f}_n^{(L^*(n))}\bigr\|_{\Lp{2}(\mu)} \\
  &\leq 2\|f\|_{\Lp{\infty}} \cdot \delta(L^*(n))^{1/2}
    + C(f,d,s) \cdot n^{-s/(2s+d)} \\
  &\leq 2\|f\|_{\Lp{\infty}} \cdot (C_0 n^{-2s/(2s+d)})^{1/2}
    + C(f,d,s) \cdot n^{-s/(2s+d)} \\
  &= O(n^{-s/(2s+d)}). \qedhere
\end{align*}
\end{proof}

\subsection{Entropy characterisation of reconstruction}
\label{III:subsec:entropy}

\begin{definition}[Fibre mixing]
\label{III:def:fibre-mixing}
Fix $\varepsilon > 0$.
An event $S^* \in \Bcal$ is \emph{$(c,\varepsilon)$-mixing across fibres} if
there exists $c > 0$ such that
\[
  \mu_z(S^*)\,\mu_z\bigl((S^*)^c\bigr) \;\geq\; c\cdot p_z
  \quad\text{for $\nuL$-a.e.\ $\varepsilon$-large fibre } z.
\]
Ergodicity of $T$ is expected to force this condition (see
Section~\ref{III:sec:discussion}), but the general derivability question remains
open; see the open problems below.
This condition is independent of~(US) and is not a standing assumption here;
it is used only in Corollary~\ref{III:cor:entropy} below.
\end{definition}

\begin{lemma}[Positive-fraction balance]
\label{III:lem:balance}
Let $S^* \in \Bcal$ be an $\varepsilon$-near-maximizer of $\deltaL$, i.e.\
$\mathbb{E}[\mathrm{Var}(\mathbbm{1}_{S^*} \mid \OhL)] \geq \deltaL - \varepsilon$.
Define the \emph{large-fibre contribution}
\[
  \kappa \;:=\; \int_{\{p_z \geq \varepsilon\}}
    p_z\,\mu_z(S^*)\,\mu_z\bigl((S^*)^c\bigr)\,d\nuL(z).
\]
For any $\eta \in (0,\tfrac{1}{2})$ with $\kappa > \eta$, the set
\[
  G(\eta) \;:=\; \bigl\{z : p_z \geq \varepsilon,\;
    \mu_z(S^*)\,\mu_z\bigl((S^*)^c\bigr) \geq \eta\bigr\}
\]
of $\varepsilon$-large balanced fibres satisfies
\[
  \nuL\bigl(G(\eta)\bigr) \;\geq\; 4(\kappa - \eta).
\]
No dynamical hypothesis is required.
\end{lemma}

\begin{proof}
Partition the $\varepsilon$-large fibres into monochromatic fibres
$\mathcal{M} = \{z : p_z \geq \varepsilon,\, \mu_z(S^*)\mu_z((S^*)^c) < \eta\}$
and balanced fibres $G(\eta)$.
On $\mathcal{M}$: $\mu_z(S^*)\mu_z((S^*)^c) < \eta$, so
\[
  \int_{\mathcal{M}} p_z\,\mu_z(S^*)\mu_z\bigl((S^*)^c\bigr)\,d\nuL
  \;\leq\; \eta\int_{\mathcal{M}} p_z\,d\nuL \;\leq\; \eta.
\]
Hence $\int_{G(\eta)} p_z\mu_z(S^*)\mu_z((S^*)^c)\,d\nuL \geq \kappa - \eta$.
On $G(\eta)$: $\mu_z(S^*)\mu_z((S^*)^c) \leq \tfrac{1}{4}$ pointwise, so
\[
  \kappa - \eta
  \;\leq\; \int_{G(\eta)} p_z\,\mu_z(S^*)\mu_z\bigl((S^*)^c\bigr)\,d\nuL
  \;\leq\; \tfrac{1}{4}\int_{G(\eta)} p_z\,d\nuL
  \;\leq\; \tfrac{1}{4}\,\nuL\bigl(G(\eta)\bigr). \qedhere
\]
\end{proof}

\begin{lemma}[Fibre mixing bound]
\label{III:lem:hard}
Let $S^*$ be a $(\deltaL - \varepsilon)$-near-maximizer of
$\mathrm{E}[\mathrm{Var}(\mathbbm{1}_{\cdot}\mid\OhL)]$ that is
$(c,\varepsilon)$-mixing across fibres
(Definition~\ref{III:def:fibre-mixing}).
Then
\[
  (\mu\otimes\mu)(\RL) \;\leq\; \frac{1}{c}\,\deltaL + C(c)\,\varepsilon,
  \qquad C(c) := \frac{1+c}{c}.
\]
\end{lemma}

\begin{proof}
Write $(\mu\otimes\mu)(\RL) = \int p_z^2\,d\nuL(z)$ and split into
$\varepsilon$-large and small-fibre parts.

\emph{Large fibres ($p_z \geq \varepsilon$).}
The $(c,\varepsilon)$-mixing condition gives
$\mu_z(S^*)\mu_z((S^*)^c) \geq c\,p_z$ on $\varepsilon$-large fibres, so:
\[
  c\int_{\{p_z \geq \varepsilon\}} p_z^2\,d\nuL(z)
  \;\leq\; \int p_z\,\mu_z(S^*)\mu_z\bigl((S^*)^c\bigr)\,d\nuL(z)
  \;=\; \mathrm{E}\bigl[\mathrm{Var}(\mathbbm{1}_{S^*}\mid\OhL)\bigr]
  \;\leq\; \deltaL.
\]
Hence $\int_{\{p_z \geq \varepsilon\}} p_z^2\,d\nuL(z) \leq \frac{1}{c}\,\deltaL$.

\emph{Small fibres ($p_z < \varepsilon$).}
Each such fibre satisfies $p_z^2 \leq \varepsilon\,p_z$, so
\[
  \int_{\{p_z < \varepsilon\}} p_z^2\,d\nuL(z)
  \;\leq\; \varepsilon\int p_z\,d\nuL(z) = \varepsilon.
\]

\emph{Combining.}
Since $C(c) = \frac{1+c}{c} \geq 1$ we have $\varepsilon \leq C(c)\,\varepsilon$,
and therefore
\[
  (\mu\otimes\mu)(\RL) \;\leq\; \frac{1}{c}\,\deltaL + C(c)\,\varepsilon. \qedhere
\]
\end{proof}

The collision probability~\eqref{eq:collision-prob} defines a natural
information-theoretic quantity:
\begin{equation}
  H_2(\nuL) \;:=\; -\log\bigl((\mu \otimes \mu)(\RL)\bigr)
  \;=\; -\log \int p_z^2\,d\nuL(z).
  \label{III:eq:entropy}
\end{equation}
This is the \emph{Rényi-$2$ (collision) entropy} of the delay-vector distribution.

\begin{corollary}[Entropy characterisation of reconstruction]
\label{III:cor:entropy}
Under the fibre mixing condition of Definition~\ref{III:def:fibre-mixing}:
\[
  \deltaL \to 0
  \quad\iff\quad
  H_2(\nuL) \to +\infty.
\]
Equivalently, reconstruction occurs if and only if the probability that two
independently drawn delay vectors coincide vanishes as $L \to \infty$.
Thus reconstruction is equivalent to the vanishing of collision probability
$(\mu \otimes \mu)(\RL)$ introduced in Section~\ref{III:subsec:fibres}.
\end{corollary}

\begin{proof}
Corollary~\ref{III:cor:easy} gives $\deltaL \leq \tfrac{1}{2}(\mu\otimes\mu)(\RL)$.
For the reverse, Lemma~\ref{III:lem:hard} with a $\deltaL$-near-maximizer $S^*$
satisfying the $(c,\varepsilon)$-mixing condition gives
$(\mu\otimes\mu)(\RL) \leq \frac{1}{c}\,\deltaL + C(c)\,\varepsilon$,
so $\deltaL \to 0 \iff (\mu\otimes\mu)(\RL) \to 0$.
Since $(\mu\otimes\mu)(\RL) = e^{-H_2(\nuL)}$, this is equivalent to
$H_2(\nuL) \to +\infty$.
\end{proof}

\begin{remark}[Balance mechanism]
\label{III:rem:balance-mechanism}
Lemma~\ref{III:lem:balance} guarantees that a positive $\nuL$-fraction of
$\varepsilon$-large fibres are $\eta$-balanced without any dynamical hypothesis.
The fibre mixing condition is the upgrade from positive-fraction to
$\nuL$-a.e.\ balance; whether ergodicity supplies it is the main open
question (Section~\ref{III:sec:discussion}).
\end{remark}

\begin{remark}[Empirical entropy witness]
\label{III:rem:entropy-witness}
The empirical analogue of $H_2(\nuL)$ is
\[
  \hat{H}_2\!\left(\nuL^{(n)}\right)
  \;:=\;
  -\log\!\left(\frac{1}{n^2}\sum_{i,j=1}^n
    \mathbbm{1}\bigl[\PhiL(x_i) = \PhiL(x_j)\bigr]\right),
\]
the negative log of the empirical collision probability.
This is the third computable witness for reconstruction.
It is \emph{computationally simpler} than $\delhat{L}{n}$: it requires no
optimisation over sets, only histogram counts over delay vectors.
Under fibre mixing, Corollary~\ref{III:cor:entropy} shows that $\Lhat$ could
equivalently be defined as the first $L$ at which $\hat{H}_2(\nuL^{(n)})$
exceeds a threshold calibrated to $n$.
\end{remark}

% ============================================================
\section{The Dynamics Theorem}
\label{III:sec:dynamics}
% ============================================================

\subsection{Lipschitz continuity of the delay map}

\begin{lemma}[Upper Lipschitz bound]
\label{III:lem:upper-lip}
Assume $T \in \mathrm{Diff}^r(X)$, $h \in C^r(X)$, $r \geq 1$.
Then $\PhiL$ is Lipschitz:
\[
  |\PhiL(x) - \PhiL(x')| \;\leq\; C_L \cdot d_X(x,x')
  \quad\forall\, x, x' \in X,
\]
where $C_L := \|Dh\|_{\Lp{\infty}} \cdot \bigl(\sum_{j=0}^{L}
\|DT^j\|_{\Lp{\infty}}^2\bigr)^{1/2}$.
\end{lemma}

\begin{proof}
Component $j$ of $\PhiL$ is $h \circ T^j$.
By the chain rule and compactness of $X$:
$|h(T^j x) - h(T^j x')| \leq \|Dh\|_{\Lp{\infty}} \cdot \|DT^j\|_{\Lp{\infty}} \cdot d_X(x,x')$.
Taking Euclidean norm over $j = 0, \ldots, L$ gives the result.
\end{proof}

\begin{lemma}[Local lower Lipschitz bound]
\label{III:lem:lower-lip-local}
Under~(G1)--(G2) with $L \geq d-1$ (Takens regime), $T \in C^r$, $r \geq 2$:
for $\mu$-a.e.\ $x \in X$, there exists a neighbourhood $U_x$ and a constant
$c_x = \sigma_{\min}(D\PhiL|_x) > 0$ such that
\[
  |\PhiL(x) - \PhiL(x')| \;\geq\; c_x \cdot d_X(x,x')
  \quad\forall\, x' \in U_x.
\]
\end{lemma}

\begin{proof}
By Lemma~\ref{III:lem:dimeff}, $D\PhiL|_x$ has full rank~$d$ for $\mu$-a.e.\ $x$
when $L \geq d-1$, so $\sigma_{\min}(D\PhiL|_x) > 0~\mu$-a.e.
The lower bound on $U_x$ follows from the inverse function theorem applied to
the immersion $\PhiL$ at $x$.
\end{proof}

\begin{remark}[Gap between local and global lower bound]
Lemma~\ref{III:lem:lower-lip-local} gives a \emph{local} lower bound $c_x > 0$
varying over $X$.
(G1)--(G2) guarantee $\sigma_{\min} > 0$~\emph{$\mu$-a.e.}, but not on all of $X$.
Condition~(US) imposes $c := \inf_X \sigma_{\min}(D\PhiL|_x) > 0$, which is
strictly stronger.
The gap is genuine: there exist pairs $(T,h)$ satisfying~(G1)--(G2) but
violating~(US) on a $\mu$-null set.
\end{remark}

\subsection{Uniform separation}

The one-step predictive kernel for deterministic $T$ is the Dirac family
$\PihL(x,\,\cdot\,) = \delta_{\PhiL(x)}$.
The TV separation pseudometric is therefore
\[
  \preddisc{L}(x,x')
  \;:=\;
  \|\PihL(x,\,\cdot\,) - \PihL(x',\,\cdot\,)\|_{\mathrm{TV}}
  \;=\;
  \mathbbm{1}\bigl[\PhiL(x) \neq \PhiL(x')\bigr]
  \;\in\; \{0,1\}.
\]
Separation is binary: $\preddisc{L}(x,x') = 1$ if and only if
$x$ and $x'$ are distinguished by $\OhL$.

\begin{theorem}[Uniform separation under~(US)]
\label{III:thm:sep-stable}
Assume~(G1)--(G2), (US), $T \in C^r$, $r \geq 2$, $L \geq d-1$.
Then $\PhiL$ is bi-Lipschitz on $X$:
\[
  c \cdot d_X(x,x') \;\leq\; |\PhiL(x) - \PhiL(x')| \;\leq\; C_L \cdot d_X(x,x')
  \quad\forall\, x, x' \in X.
\]
For any estimator $\GamHat$ with uniform error
$|\GamHat(x) - \PhiL(x)| \leq \varepsilon_n$ uniformly in $x$, the
empirical separation $\widehat{d}_L(x,x') := |\GamHat(x) - \GamHat(x')|$
satisfies
\[
  \widehat{d}_L(x,x') \;\geq\; c \cdot d_X(x,x') - 2\varepsilon_n.
\]
In particular, if $\preddisc{L}(x,x') = 1$ (i.e.\ $\PhiL(x) \neq \PhiL(x')$),
then $\widehat{d}_L(x,x') > 0$ for all $n$ large enough that
$\varepsilon_n < c \cdot d_X(x,x') / 2$.
\end{theorem}

\begin{proof}
Upper bound: Lemma~\ref{III:lem:upper-lip}.
Lower bound: (US) directly gives $|\PhiL(x) - \PhiL(x')| \geq c \cdot d_X(x,x')$
for all $x, x' \in X$.
For the empirical bound: by the triangle inequality,
$|\GamHat(x) - \GamHat(x')| \geq |\PhiL(x) - \PhiL(x')| - 2\varepsilon_n
\geq c \cdot d_X(x,x') - 2\varepsilon_n$.
If $\PhiL(x) \neq \PhiL(x')$ then by bi-Lipschitz $d_X(x,x') \geq |\PhiL(x)-\PhiL(x')|/C_L > 0$,
so the right side is eventually positive.
\end{proof}

\begin{theorem}[Dynamics Theorem]
\label{III:thm:dynamics}
Assume~(G1)--(G2), (US), $T \in C^r$, $r \geq 2$, $\mu$ smooth positive
density on compact $d$-manifold $X$, $L \geq d-1$.
Let $\GamHat$ be the empirical delay map estimator with uniform error
$\varepsilon_n = C'' n^{-\beta/(2\beta+d)}$ (Proposition~\ref{III:prop:source2}).
Then:
\begin{enumerate}[label=(\alph*)]
  \item \emph{(Convergence)} $\|\GamHat - \PhiL\|_{\Lp{\infty}} \to 0$ at
    rate $n^{-\beta/(2\beta+d)}$.
  \item \emph{(Separation)} For $\mu \otimes \mu$-a.e.\ $(x,x')$ with
    $x \neq x'$: $\widehat{d}_L(x,x') > 0$ for all large $n$.
  \item \emph{(Scope)} Condition~(US) is not a consequence of~(G1)--(G2).
    Without~(US), the kernel witness $\widehat{d}_L$ can vanish on a positive-measure
    set of pairs even when $\preddisc{L} = 1$ and reconstruction holds.
\end{enumerate}
\end{theorem}

\begin{proof}
(a) Proposition~\ref{III:prop:source2} with $\deff(L) = d$ gives
$\varepsilon_n \to 0$ at rate $n^{-\beta/(2\beta+d)}$.

(b) Since $\mu$ is non-atomic (smooth positive density), $\mu \otimes \mu$
almost every pair satisfies $d_X(x,x') > 0$.
For such pairs $\preddisc{L}(x,x') = 1$ (reconstruction and injectivity of $\PhiL$),
and by Theorem~\ref{III:thm:sep-stable}, $\widehat{d}_L(x,x') \geq c \cdot d_X(x,x') - 2\varepsilon_n$.
For fixed $x \neq x'$, $c \cdot d_X(x,x') > 0$ and $\varepsilon_n \to 0$,
so the right side is eventually positive.

(c) The proof of Theorem~\ref{III:thm:sep-stable} uses~(US) in the lower bound.
Without it, there exist points $x$ where $\sigma_{\min}(D\PhiL|_x) \to 0$
along sequences $x'_n \to x$, and $|\PhiL(x_n') - \PhiL(x)| / d_X(x,x'_n) \to 0$.
Any estimator inherits this degeneracy.
\end{proof}

% ============================================================
\section{The Conjunction Theorem}
\label{III:sec:conjunction}
% ============================================================

This section formalises the complete reconstruction diagnostic and proves
that both witnesses certify the same event when reconstruction and~(US) hold.

\subsection{The honest bridge}

The delay vector $\PhiL(x) = (h(x), \ldots, h(T^L x))$ is a trajectory of the
Markov chain on $h$-space with one-step kernel $\Pi_h$.
The $\sigma$-algebra $\OhL = \sigma(\PhiL)$ is exactly the information
generated by $L$ steps of this Markov chain starting from $x$.
This gives an algebraic identity:
\[
  \sigma\bigl(\PhiL\bigr) \;=\; \OhL,
\]
which is not an approximation but a definition.
Consequently, the $\sigma$-algebra separation event and the metric separation event
are identical:

\begin{lemma}[Algebra separation equals metric separation]
\label{III:lem:alg-metric}
For deterministic $T$ and $\mu \otimes \mu$-a.e.\ $(x, x')$:
\[
  \preddisc{L}(x,x') = 1
  \;\iff\;
  \text{$x$ and $x'$ are separated by $\OhL$.}
\]
These are the same event, not merely correlated.
\end{lemma}

\begin{proof}
For deterministic $T$:
$\PihL(x,\,\cdot\,) = \delta_{\PhiL(x)}$, so
$\preddisc{L}(x,x') = \|\delta_{\PhiL(x)} - \delta_{\PhiL(x')}\|_{\mathrm{TV}}
= \mathbbm{1}[\PhiL(x) \neq \PhiL(x')]$.
And $x$ is separated from $x'$ by $\OhL = \sigma(\PhiL)$ if and only if
$\exists E \in \OhL$ with $\mathbbm{1}_E(x) \neq \mathbbm{1}_E(x')$, i.e.,
if and only if $\PhiL(x) \neq \PhiL(x')$ (since $\OhL$ is generated by the
level sets of $\PhiL$).
The two conditions are identical.
\end{proof}

Under reconstruction ($\deltaL = 0$, equivalently $\PhiL$ injective $\mu$-a.e.)
and~(US), injectivity holds everywhere, so $\preddisc{L}(x,x') = 1$ for
$\mu \otimes \mu$-a.e.\ $x \neq x'$.

\subsection{The joint diagnostic}

\begin{definition}[Joint reconstruction diagnostic]
\label{III:def:diagnostic}
Fix tolerances $\tau_n, \eta_n > 0$.
Define
\[
  \mathrm{Reconstruct}(L, n)
  \;:=\;
  \bigl\{\delhat{L}{n} \leq \tau_n\bigr\}
  \;\wedge\;
  \bigl\{\widehat{d}_L(x,x') > \eta_n
    \text{ for $\mu \otimes \mu$-most pairs}\bigr\}.
\]
\end{definition}

\begin{theorem}[Conjunction Theorem]
\label{III:thm:conjunction}
Assume the hypotheses of Theorem~\ref{III:thm:algebra} (for the algebra side)
and Theorem~\ref{III:thm:dynamics} (for the dynamics side).
Set $\tau_n = \frac{C_\lambda}{2} C' n^{-2s/(2s+d)}$ and
$\eta_n = \sqrt{\varepsilon_n}$ where $\varepsilon_n = C' n^{-2s/(2s+d)}$.
Then:
\begin{enumerate}[label=(\alph*)]
  \item \emph{(Conjunction fires under reconstruction~$+$~(US)).}
    Under $\deltaL = 0$ for $L \geq L^*(n)$ and~(US):
    \[
      P\bigl(\mathrm{Reconstruct}(\Lhat, n)\bigr) \;\to\; 1
      \quad\text{exponentially in } n.
    \]
  \item \emph{(Algebra witness fails when reconstruction fails).}
    If $\delta(L) > 0$ for all $L$: $P(\delhat{L}{n} \leq \tau_n) \to 0$
    for every fixed $L$, and $\Lhat \to \infty$.

  \item \emph{(Kernel witness fails when~(US) fails).}
    If $\sigma_{\min}(D\PhiL|_x) \to 0$ on a positive-measure set $S$:
    $P(\widehat{d}_L > \eta_n \text{ for }\mu\otimes\mu\text{-most pairs}) \to 0$
    even when $\delhat{L}{n} \to 0$.
\end{enumerate}
In particular: the conjunction is not achievable when either hypothesis fails.
\end{theorem}

\begin{proof}
\emph{Part~(a).}
By Theorem~\ref{III:thm:elbow}, $P(\Lhat = L^*(n)) \to 1$ exponentially.
On $\{\Lhat = L^*(n)\}$:
(i) $\delhat{\Lhat}{n} \leq \tau_n$ whp by Theorem~\ref{III:thm:convergence}
and the definition of $L^*(n)$ ($\delta(L^*(n)) = 0$ under reconstruction).
(ii) By Theorem~\ref{III:thm:dynamics}(b), $\widehat{d}_{\Lhat}(x,x') > 0$
for $\mu \otimes \mu$-a.e.\ $(x,x')$ with $x \neq x'$, eventually.
The set $\{d_X(x,x') \leq (2\varepsilon_n + \eta_n)/c\}$ has
$\mu \otimes \mu$-measure $\to 0$ as $\eta_n/c \to 0$, so
$\widehat{d}_{\Lhat} > \eta_n$ for $\mu\otimes\mu$-most pairs.

\emph{Part~(b).}
If $\delta(L) > 0$ for all $L$, then by Theorem~\ref{III:thm:convergence},
$\delhat{L}{n} \geq \delta(L) - \varepsilon_n \geq \delta(L)/2 > \tau_n$
with high probability for all fixed $L$, so the first condition of
$\mathrm{Reconstruct}$ never fires.

\emph{Part~(c).}
If $S = \{x : \sigma_{\min}(D\PhiL|_x) = 0\}$ has $\mu(S) > 0$,
then for $x \in S$ there exist sequences $x'_k \to x$ with
$|\PhiL(x'_k) - \PhiL(x)| / d_X(x'_k, x) \to 0$.
Any consistent estimator satisfies
$\widehat{d}_L(x, x'_k) \leq |\PhiL(x'_k) - \PhiL(x)| + 2\varepsilon_n \to 0$.
The set of pairs where $\widehat{d}_L \leq \eta_n$ thus has positive
$\mu \otimes \mu$-measure, and the kernel condition fails.
\end{proof}

% ============================================================
\section{Discussion}
\label{III:sec:discussion}
% ============================================================

The results of this paper divide into two layers.  Section~\ref{III:sec:general}
establishes general observational certification: the population separation defect
$\delta(\mathcal{G})$ converges to zero under CE and honest refinement, the
empirical U-statistic proxy concentrates under IID sampling, and a persistent
empirical floor is evidence against CE.  These results require no temporal or
geometric structure.  Sections~\ref{III:sec:algebra}--\ref{III:sec:conjunction}
establish the dynamical sharpening: honest refinement is automatic under lag
growth, and mixing upgrades consistency to minimax-optimal rates with computable
witnesses.  The open problems below all concern the dynamical layer; the general
layer is complete at the level of consistency.

\textbf{Polynomial mixing.}
The Algebra Theorem (Theorem~\ref{III:thm:algebra}) and the Elbow Location Theorem
(Theorem~\ref{III:thm:elbow}) require exponential mixing with resolvability constant
$C_\lambda > 4$.
Under polynomial mixing $\deltaL \lesssim L^{-\alpha}$, the elbow drop at
$L^*(n)$ satisfies $\Delta(L^*(n)) \sim \varepsilon_n / L^*(n) \to 0$ relative
to $\varepsilon_n$: the elbow and the noise floor merge asymptotically and
cannot be separated.
A different stopping criterion (e.g.\ integrated elbow, or a bandwidth-adaptive
rule) is needed.

\textbf{The $d > 2s$ hypothesis.}
Proposition~\ref{III:prop:concentration} requires $d > 2s$ for the Rademacher
concentration to be useful at rate $\varepsilon_n$.
When $d \leq 2s$, localised Rademacher complexity
(Bartlett--Bousquet--Mendelson~\cite{BartlettBousquetMendelson2005})
gives adaptive concentration without this condition but requires additional
work.

\textbf{Uniform separation from first principles.}
Condition~(US) is a hypothesis in Theorems~\ref{III:thm:dynamics}
and~\ref{III:thm:conjunction}.
The natural conjecture is that $T$ Anosov with $h$ generic implies~(US) via
the uniform hyperbolicity constants.
This would close the gap between~(G1)--(G2) and~(US) in the hyperbolic category.

\textbf{Stochastic $T$.}
For stochastic dynamics, $\PihL(x,\,\cdot\,)$ is not a Dirac delta.
The TV separation $\preddisc{L}$ is no longer binary, and Hölder continuity
of $x \mapsto \PihL(x,\,\cdot\,)$ in TV is both necessary and non-trivial.
The four-link chain of Proposition~\ref{III:prop:source2} adapts, but the
Conjunction Theorem's proof structure changes substantially.

\textbf{Sharp rates for the data-driven lag.}
Theorem~\ref{III:thm:elbow} shows $P(\Lhat = L^*(n)) \to 1$ but does not bound
the fluctuations of $\Lhat$ around $L^*(n)$ on the event $\{\Lhat \neq L^*(n)\}$.
A Berry--Esseen-type result for $\Lhat - L^*(n)$ would sharpen the stopping rule.

\textbf{Concentration of the entropy witness.}
The empirical collision entropy $\hat{H}_2(\nuL^{(n)})$ is a computationally
simpler stopping criterion (Remark~\ref{III:rem:entropy-witness}).
Asymptotic equivalence with $\delhat{L}{n}$ is established under fibre mixing
(Corollary~\ref{III:cor:entropy}), but a finite-$n$ concentration result for
$\hat{H}_2(\nuL^{(n)})$ --- comparable to Proposition~\ref{III:prop:concentration}
for $\delhat{L}{n}$ --- is not proved here.
Such a result requires only McDiarmid-type bounds on the U-statistic
$\frac{1}{n^2}\sum_{i,j} \mathbbm{1}[\PhiL(x_i) = \PhiL(x_j)]$, with no
new structural assumptions beyond those already made.

\textbf{Derivability of fibre mixing (the main conceptual open problem).}
The central open question is whether the fibre mixing condition
(Definition~\ref{III:def:fibre-mixing}) is derivable from dynamical hypotheses, or
whether it is an irreducible admissibility condition analogous to CE in~\cite{mahon_paper1}.

Lemma~\ref{III:lem:balance} settles Step~A: a positive $\nuL$-fraction of
$\varepsilon$-large fibres are $\eta$-balanced, with an explicit lower bound
depending only on the large-fibre contribution $\kappa$ and not on any dynamical
hypothesis.

The open question is Step~B: what is the weakest dynamical condition that
upgrades positive-fraction balance to $\nuL$-a.e.\ balance?  Ergodicity is the
natural candidate: in an ergodic system, a near-maximizer of $\deltaL$ must
cross fibres rather than align with them, since aligned events are
$\mathcal{O}_h^{(L)}$-measurable and thus carry zero approximation error.  The
informal argument is clear, but formalizing the connection between ergodicity
and the a.e.\ balance condition is the isolated open gap.  If ergodicity
suffices, the fibre mixing condition is derivable and the structural analogy
with CE breaks; if it does not, the analogy is complete.

\ifstandalone
\bibliographystyle{plain}
\bibliography{references}
\fi
